\batchmode
\documentclass[12pt]{article}
\makeatletter
\usepackage{epsf,amsmath}

\oddsidemargin0in
\textwidth6.5in
\topmargin-0.5in
\textheight9in
\parindent0in

\providecommand{\OneFigure}[1]{\begin{center}
\parbox{1.5in}{\epsfysize=1.5in \epsfbox{#1}}\end{center}}
\par\usepackage[dvips]{color}
\pagecolor[gray]{.7}


\makeatletter
\count@=\the\catcode`\_ \catcode`\_=8 
\newenvironment{tex2html_wrap}{}{} \catcode`\_=\count@
\makeatother
\let\mathon=$
\let\mathoff=$
\ifx\AtBeginDocument\undefined \newcommand{\AtBeginDocument}[1]{}\fi
\newbox\sizebox
\setlength{\hoffset}{0pt}\setlength{\voffset}{0pt}
\addtolength{\textheight}{\footskip}\setlength{\footskip}{0pt}
\addtolength{\textheight}{\topmargin}\setlength{\topmargin}{0pt}
\addtolength{\textheight}{\headheight}\setlength{\headheight}{0pt}
\addtolength{\textheight}{\headsep}\setlength{\headsep}{0pt}
\setlength{\textwidth}{349pt}
\newwrite\lthtmlwrite
\makeatletter
\let\realnormalsize=\normalsize
\global\topskip=2sp
\def\preveqno{}\let\real@float=\@float \let\realend@float=\end@float
\def\@float{\let\@savefreelist\@freelist\real@float}
\def\end@float{\realend@float\global\let\@freelist\@savefreelist}
\let\real@dbflt=\@dbflt \let\end@dblfloat=\end@float
\let\@largefloatcheck=\relax
\def\@dbflt{\let\@savefreelist\@freelist\real@dbflt}
\def\adjustnormalsize{\def\normalsize{\mathsurround=0pt \realnormalsize
 \parindent=0pt\abovedisplayskip=0pt\belowdisplayskip=0pt}\normalsize}%
\def\lthtmltypeout#1{{\let\protect\string\immediate\write\lthtmlwrite{#1}}}%
\newcommand\lthtmlhboxmathA{\adjustnormalsize\setbox\sizebox=\hbox\bgroup}%
\newcommand\lthtmlvboxmathA{\adjustnormalsize\setbox\sizebox=\vbox\bgroup%
 \let\ifinner=\iffalse }%
\newcommand\lthtmlboxmathZ{\@next\next\@currlist{}{\def\next{\voidb@x}}%
 \expandafter\box\next\egroup}%
\newcommand\lthtmlmathtype[1]{\def\lthtmlmathenv{#1}}%
\newcommand\lthtmllogmath{\lthtmltypeout{l2hSize %
:\lthtmlmathenv:\the\ht\sizebox::\the\dp\sizebox::\the\wd\sizebox.\preveqno}}%
\newcommand\lthtmlfigureA[1]{\let\@savefreelist\@freelist
       \lthtmlmathtype{#1}\lthtmlvboxmathA}%
\newcommand\lthtmlfigureZ{\lthtmlboxmathZ\lthtmllogmath\copy\sizebox
       \global\let\@freelist\@savefreelist}%
\newcommand\lthtmldisplayA[1]{\lthtmlmathtype{#1}\lthtmlvboxmathA}%
\newcommand\lthtmldisplayB[1]{\edef\preveqno{(\theequation)}%
  \lthtmldisplayA{#1}\let\@eqnnum\relax}%
\newcommand\lthtmldisplayZ{\lthtmlboxmathZ\lthtmllogmath\lthtmlsetmath}%
\newcommand\lthtmlinlinemathA[1]{\lthtmlmathtype{#1}\lthtmlhboxmathA  \vrule height1.5ex width0pt }%
\newcommand\lthtmlinlineA[1]{\lthtmlmathtype{#1}\lthtmlhboxmathA}%
\newcommand\lthtmlinlineZ{\egroup\expandafter\ifdim\dp\sizebox>0pt %
  \expandafter\centerinlinemath\fi\lthtmllogmath\lthtmlsetinline}
\newcommand\lthtmlinlinemathZ{\egroup\expandafter\ifdim\dp\sizebox>0pt %
  \expandafter\centerinlinemath\fi\lthtmllogmath\lthtmlsetmath}
\def\lthtmlsetinline{\hbox{\vrule width.1em\vtop{\vbox{%
  \kern.1em\copy\sizebox}\ifdim\dp\sizebox>0pt\kern.1em\else\kern.3pt\fi
  \ifdim\hsize>\wd\sizebox \hrule depth1pt\fi}}}
\def\lthtmlsetmath{\hbox{\vrule width.1em\vtop{\vbox{%
  \kern.1em\kern0.8 pt\hbox{\hglue.17em\copy\sizebox\hglue0.8 pt}}\kern.3pt%
  \ifdim\dp\sizebox>0pt\kern.1em\fi \kern0.8 pt%
  \ifdim\hsize>\wd\sizebox \hrule depth1pt\fi}}}
\def\centerinlinemath{%\dimen1=\ht\sizebox
  \dimen1=\ifdim\ht\sizebox<\dp\sizebox \dp\sizebox\else\ht\sizebox\fi
  \advance\dimen1by.5pt \vrule width0pt height\dimen1 depth\dimen1 
 \dp\sizebox=\dimen1\ht\sizebox=\dimen1\relax}

\def\lthtmlcheckvsize{\ifdim\ht\sizebox<\vsize\expandafter\vfill
  \else\expandafter\vss\fi}%
\makeatletter \tracingstats = 1 


\begin{document}
\pagestyle{empty}\thispagestyle{empty}%
\lthtmltypeout{latex2htmlLength hsize=\the\hsize}%
\lthtmltypeout{latex2htmlLength vsize=\the\vsize}%
\lthtmltypeout{latex2htmlLength hoffset=\the\hoffset}%
\lthtmltypeout{latex2htmlLength voffset=\the\voffset}%
\lthtmltypeout{latex2htmlLength topmargin=\the\topmargin}%
\lthtmltypeout{latex2htmlLength topskip=\the\topskip}%
\lthtmltypeout{latex2htmlLength headheight=\the\headheight}%
\lthtmltypeout{latex2htmlLength headsep=\the\headsep}%
\lthtmltypeout{latex2htmlLength parskip=\the\parskip}%
\lthtmltypeout{latex2htmlLength oddsidemargin=\the\oddsidemargin}%
\makeatletter
\if@twoside\lthtmltypeout{latex2htmlLength evensidemargin=\the\evensidemargin}%
\else\lthtmltypeout{latex2htmlLength evensidemargin=\the\oddsidemargin}\fi%
\makeatother
{\newpage\clearpage
\lthtmldisplayA{displaymath242}%
\begin{displaymath}\int_0^{\pi^2} \frac{t^3-3e^t}{5-\cos(\sqrt{t})} dt, \ \ \ \ \ 
\int_0^{\pi^2} \frac{t^3-3e^t}{5-\cos(t)} dt, \ \ \ \ \ 
\int_0^{\pi^2} \frac{t^3-3e^t}{1-\cos(\sqrt{t})} dt, \ \ \ \ \ 
\int_0^{\pi^2} \frac{t^3-3e^t}{1-\sin(t)} dt.
\end{displaymath}%
\lthtmldisplayZ
\hfill\lthtmlcheckvsize\clearpage}

{\newpage\clearpage
\lthtmlinlinemathA{tex2html_wrap_inline288}%
$t=4\pi$%
\lthtmlinlinemathZ
\hfill\lthtmlcheckvsize\clearpage}

{\newpage\clearpage
\lthtmlinlinemathA{tex2html_wrap_inline290}%
$y(4\pi)$%
\lthtmlinlinemathZ
\hfill\lthtmlcheckvsize\clearpage}

{\newpage\clearpage
\lthtmlinlinemathA{tex2html_wrap_inline292}%
$\frac{\pi}{16}$%
\lthtmlinlinemathZ
\hfill\lthtmlcheckvsize\clearpage}

{\newpage\clearpage
\lthtmlinlinemathA{tex2html_wrap_inline294}%
$\frac{\pi}{32}$%
\lthtmlinlinemathZ
\hfill\lthtmlcheckvsize\clearpage}

{\newpage\clearpage
\lthtmlinlinemathA{tex2html_wrap_inline296}%
$\frac{\pi}{64}$%
\lthtmlinlinemathZ
\hfill\lthtmlcheckvsize\clearpage}

{\newpage\clearpage
\lthtmldisplayA{eqnarraystar46}%
\begin{eqnarray*}|I_{\frac{\pi}{16}}-I_{\frac{\pi}{32}}| & = & 0.0340391388, \\
|I_{\frac{\pi}{32}}-I_{\frac{\pi}{64}}| & = & 0.0192935453.
\end{eqnarray*}%
\lthtmldisplayZ
\hfill\lthtmlcheckvsize\clearpage}

{\newpage\clearpage
\lthtmlinlinemathA{tex2html_wrap_inline308}%
$\pi$%
\lthtmlinlinemathZ
\hfill\lthtmlcheckvsize\clearpage}

{\newpage\clearpage
\lthtmlinlinemathA{tex2html_wrap_inline312}%
$\frac{\pi}{1024}$%
\lthtmlinlinemathZ
\hfill\lthtmlcheckvsize\clearpage}

{\newpage\clearpage
\lthtmlinlinemathA{tex2html_wrap_inline314}%
$\int_0^1 f(x) dx$%
\lthtmlinlinemathZ
\hfill\lthtmlcheckvsize\clearpage}

{\newpage\clearpage
\lthtmlinlinemathA{tex2html_wrap_inline318}%
$f'\left(\frac{1}{4}\right)$%
\lthtmlinlinemathZ
\hfill\lthtmlcheckvsize\clearpage}

{\newpage\clearpage
\lthtmldisplayA{displaymath246}%
\begin{displaymath}\int_0^1 f(x) dx \approx A_1 f(0) + A_2 f'\left(\frac{1}{4}\right)+
A_3 f(1). \ \ 
\footnote{While it is not relevant to this exam, 
if you were using this over a general panel, it would have
the form
\begin{displaymath}
\int_{x_i}^{x_{i+1}} f(x) dx \approx h \left[A_1 f(x_i) + 
h A_2 f'\left(x_{i+\frac{1}{4}}\right)+
A_3 f(x_{i+1})\right].
\end{displaymath}
}
\end{displaymath}%
\lthtmldisplayZ
\hfill\lthtmlcheckvsize\clearpage}

{\newpage\clearpage
\lthtmlinlinemathA{tex2html_wrap_inline328}%
$\xi = 1/2$%
\lthtmlinlinemathZ
\hfill\lthtmlcheckvsize\clearpage}

{\newpage\clearpage
\lthtmldisplayA{eqnarraystar68}%
\begin{eqnarray*}c_0 & = & f(0), \\
c_1 & = & f(0)-f(1) + 2 f'\left(\frac{1}{4}\right), \\
c_2 & = & 2 f(1) - 2 f(0) - 2 f'\left(\frac{1}{4}\right).
\end{eqnarray*}%
\lthtmldisplayZ
\hfill\lthtmlcheckvsize\clearpage}

{\newpage\clearpage
\lthtmldisplayA{displaymath249}%
\begin{displaymath}\int_0^1 p(x) dx = A_1 f(0) + A_2 f'\left(\frac{1}{4}\right) + A_3 f(1),
\end{displaymath}%
\lthtmldisplayZ
\hfill\lthtmlcheckvsize\clearpage}

{\newpage\clearpage
\lthtmldisplayA{eqnarraystar76}%
\begin{eqnarray*}A_1 & = & \frac{5}{6}, \\
A_2 & = & \frac{1}{3}, \\
A_3 & = & \frac{1}{6}.
\end{eqnarray*}%
\lthtmldisplayZ
\hfill\lthtmlcheckvsize\clearpage}

{\newpage\clearpage
\lthtmldisplayA{displaymath250}%
\begin{displaymath}\int_0^1 f(x) dx = 
\frac{5}{6} f(0) + \frac{1}{3} f'\left(\frac{1}{4}\right)+
\frac{1}{6} f(1) + O(h^3).\end{displaymath}%
\lthtmldisplayZ
\hfill\lthtmlcheckvsize\clearpage}

{\newpage\clearpage
\lthtmldisplayA{displaymath251}%
\begin{displaymath}y' = f(y,t), \ \ \ \ y(0) = y_0.
\end{displaymath}%
\lthtmldisplayZ
\hfill\lthtmlcheckvsize\clearpage}

{\newpage\clearpage
\lthtmldisplayA{displaymath252}%
\begin{displaymath}u_{k+1} = u_k + h f(u_k,t_k), \ \ \ \ u_0 = y_0.
\end{displaymath}%
\lthtmldisplayZ
\hfill\lthtmlcheckvsize\clearpage}

{\newpage\clearpage
\lthtmldisplayA{displaymath253}%
\begin{displaymath}y'(t_k) = f(u_k,t_k) \approx \frac{u_k - u_{k-1}}{h}
\end{displaymath}%
\lthtmldisplayZ
\hfill\lthtmlcheckvsize\clearpage}

{\newpage\clearpage
\lthtmldisplayA{displaymath254}%
\begin{displaymath}u_{k+1} = u_k + h \left( \frac{u_k - u_{k-1}}{h}\right), \ \ \ \ u_0 = y_0.
\end{displaymath}%
\lthtmldisplayZ
\hfill\lthtmlcheckvsize\clearpage}

{\newpage\clearpage
\lthtmldisplayA{eqnarraystar106}%
\begin{eqnarray*}y_{k+1} & = & y_k + y'_k h + \frac{1}{2} y''_k h^2 + \frac{1}{3!}
y''' (\xi_1) h^3 \\
y_{k+1} & = & y_{k-1} + y'_{k-1} (2h) + \frac{1}{2} y''_{k-1} (2h)^2 +
\frac{1}{3!} y'''(\xi_2) (2h)^3
\end{eqnarray*}%
\lthtmldisplayZ
\hfill\lthtmlcheckvsize\clearpage}

{\newpage\clearpage
\lthtmldisplayA{eqnarraystar121}%
\begin{eqnarray*}y_{k+1}-u_{k+1} & = & y_{k+1}-2 u_k + u_{k-1} \\
& = & y_{k+1}-2 y_k + y_{k-1} \\
& = & 2 y_{k+1} - y_{k+1} -2 y_k + y_{k-1} \\
& = & 2 \left[ y'_k h + \frac{1}{2} y''_k h^2 + \frac{1}{3!}
y'''(\xi_1) h^3\right] - \left[y'_{k-1} (2h) + \frac{1}{2} y''_{k-1} (2h)^2 +
\frac{1}{3!} y'''(\xi_2) (2h)^3\right] \\
& = & 2 h \left[ y'_k - y'_{k-1} \right]+ 
h^2 \left[y''_k - 2 y''_{k-1}\right] + \frac{1}{3!}
y'''(\xi_1) h^3 -
\frac{1}{3!} y'''(\xi_2) (2h)^3 \\
& = & 2 h \left[ y'_k - y'_k \right]+ 
h^2 \left[y''_k - 2 y''_{k} + 2 y''_k \right] + \\
& & \frac{1}{3!}
y'''(\xi_1) h^3 -
\frac{1}{3!} y'''(\xi_2) (2h)^3 +  y'''(\xi_3) h^3 +  \frac{1}{2}
y'''(\xi_4) h^3 \\
& = &  h^2 y''_k + \frac{1}{3!}
y'''(\xi_1) h^3 -
\frac{1}{3!} y'''(\xi_2) (2h)^3 +  y'''(\xi_3) h^3 +  \frac{1}{2}
y'''(\xi_4) h^3 \\
& = & O(h^2)
\end{eqnarray*}%
\lthtmldisplayZ
\hfill\lthtmlcheckvsize\clearpage}

{\newpage\clearpage
\lthtmldisplayA{displaymath257}%
\begin{displaymath}y_{k+1} - u_{k+1} = \frac{1}{2} f''(\xi) h^2 = h\tau_k = O(h^2)
\end{displaymath}%
\lthtmldisplayZ
\hfill\lthtmlcheckvsize\clearpage}

{\newpage\clearpage
\lthtmldisplayA{displaymath258}%
\begin{displaymath}e_{n+1} = e_n + h[f(y_{n+1},t_{n+1})-f(u_{n+1},t_{n+1})] \\
\end{displaymath}%
\lthtmldisplayZ
\hfill\lthtmlcheckvsize\clearpage}

{\newpage\clearpage
\lthtmldisplayA{eqnarraystar171}%
\begin{eqnarray*}(1-hK) |e_{N+1}| & \leq & |e_N| + h \tau \\
|e_{N+1}| & \leq & \frac{1}{1-hK}|e_N| + \frac{1}{1-hK} h \tau
\end{eqnarray*}%
\lthtmldisplayZ
\hfill\lthtmlcheckvsize\clearpage}

{\newpage\clearpage
\lthtmldisplayA{eqnarraystar179}%
\begin{eqnarray*}|e_{N+1}| & \leq & \frac{1}{(1-hK)^{T/h}} | e_0| +
\frac{(1-hK)^{-N}-1}{(1-hK)^{-1}-1} h \tau \\
& \leq & e^{KT} | e_0| +
\frac{e^{KT}-1}{K} \tau.
\end{eqnarray*}%
\lthtmldisplayZ
\hfill\lthtmlcheckvsize\clearpage}

{\newpage\clearpage
\lthtmldisplayA{displaymath259}%
\begin{displaymath}u_{k+1} = u_k + h\left[\frac{2}{3} f(u_k,t) + \frac{1}{3} f(u_{k+1},t)\right].
\end{displaymath}%
\lthtmldisplayZ
\hfill\lthtmlcheckvsize\clearpage}

{\newpage\clearpage
\lthtmldisplayA{displaymath260}%
\begin{displaymath}y' = -\lambda y,
\end{displaymath}%
\lthtmldisplayZ
\hfill\lthtmlcheckvsize\clearpage}

{\newpage\clearpage
\lthtmlinlinemathA{tex2html_wrap_inline358}%
$\lambda>0$%
\lthtmlinlinemathZ
\hfill\lthtmlcheckvsize\clearpage}

{\newpage\clearpage
\lthtmlinlinemathA{tex2html_wrap_inline360}%
$h\lambda$%
\lthtmlinlinemathZ
\hfill\lthtmlcheckvsize\clearpage}

{\newpage\clearpage
\lthtmlinlinemathA{tex2html_wrap_inline364}%
$y' = -\lambda y$%
\lthtmlinlinemathZ
\hfill\lthtmlcheckvsize\clearpage}

{\newpage\clearpage
\lthtmldisplayA{displaymath261}%
\begin{displaymath}u_{k+1} = u_k - \frac{2}{3} h \lambda u_k - \frac{1}{3} h \lambda u_{k+1}
\end{displaymath}%
\lthtmldisplayZ
\hfill\lthtmlcheckvsize\clearpage}

{\newpage\clearpage
\lthtmldisplayA{displaymath262}%
\begin{displaymath}u_{k+1} = \left(\frac{1-\frac{2}{3}h\lambda}{1+\frac{1}{3}h\lambda}\right)u_k.
\end{displaymath}%
\lthtmldisplayZ
\hfill\lthtmlcheckvsize\clearpage}

{\newpage\clearpage
\lthtmldisplayA{displaymath263}%
\begin{displaymath}u_k = \left(\frac{1-\frac{2}{3}h\lambda}{1+\frac{1}{3}h\lambda}\right)^k,
\end{displaymath}%
\lthtmldisplayZ
\hfill\lthtmlcheckvsize\clearpage}

{\newpage\clearpage
\lthtmlinlinemathA{tex2html_wrap_inline368}%
$h\lambda>3/2$%
\lthtmlinlinemathZ
\hfill\lthtmlcheckvsize\clearpage}

{\newpage\clearpage
\lthtmlinlinemathA{tex2html_wrap_inline370}%
$h \lambda > 6$%
\lthtmlinlinemathZ
\hfill\lthtmlcheckvsize\clearpage}


\end{document}
